% problem_id: S001-lemma-complexes-open-neg-exts-vanishing
% source_id: S001 (Stacks Project)
% source_file: quot.tex
% statement_locator: quot.tex lines 4936--4965
% proof_locator: quot.tex lines 4967--5102
% env: lemma
% id_note: label 在本来源内唯一
% pairing: tight (verified)
% status: complete (statement + author proof verbatim + reason + locators)
%
\begin{lemma}
\label{lemma-complexes-open-neg-exts-vanishing}
Let $S$ be a scheme.
Let $f : X \to Y$ be a morphism of algebraic spaces over $S$.
Assume $f$ is proper, flat, and of finite presentation.
Let $K, E \in D(\mathcal{O}_X)$. Assume $K$ is pseudo-coherent
and $E$ is $Y$-perfect (More on Morphisms of Spaces, Definition
\ref{spaces-more-morphisms-definition-relatively-perfect}).
For a field $k$ and a morphism $y : \Spec(k) \to Y$ denote $K_y$, $E_y$
the pullback to the fibre $X_y$.
\begin{enumerate}
\item There is an open $W \subset Y$ characterized by the property
$$
y \in |W|
\Leftrightarrow
\Ext^i_{\mathcal{O}_{X_y}}(K_y, E_y) = 0
\text{ for }i < 0.
$$
\item For any morphism $V \to Y$ factoring through $W$ we have
$$
\Ext^i_{\mathcal{O}_{X_V}}(K_V, E_V) = 0
\quad\text{for}\quad i < 0
$$
where $X_V$ is the base change of $X$ and $K_V$ and $E_V$
are the derived pullbacks of $K$ and $E$ to $X_V$.
\item The functor $V \mapsto \Hom_{\mathcal{O}_{X_V}}(K_V, E_V)$
is a sheaf on $(\textit{Spaces}/W)_{fppf}$ representable by an
algebraic space affine and of finite presentation over $W$.
\end{enumerate}
\end{lemma}

\begin{proof}
For any morphism $V \to Y$ the complex $K_V$ is pseudo-coherent
(Cohomology on Sites, Lemma
\ref{sites-cohomology-lemma-pseudo-coherent-pullback})
and $E_V$ is $V$-perfect (More on Morphisms of Spaces, Lemma
\ref{spaces-more-morphisms-lemma-base-change-relatively-perfect}).
Another observation is that given $y : \Spec(k) \to Y$
and a field extension $k'/k$ with $y' : \Spec(k') \to Y$
the induced morphism, we have
$$
\Ext^i_{\mathcal{O}_{X_{y'}}}(K_{y'}, E_{y'}) =
\Ext^i_{\mathcal{O}_{X_y}}(K_y, E_y) \otimes_k k'
$$
by Derived Categories of Schemes, Lemma
\ref{perfect-lemma-affine-morphism-and-hom-out-of-perfect}.
Thus the vanishing in (1) is really a property of the induced
point $y \in |Y|$.
We will use these two observations without further mention in the proof.

\medskip\noindent
Assume first $Y$ is an affine scheme. Then we may apply
More on Morphisms of Spaces, Lemma
\ref{spaces-more-morphisms-lemma-compute-ext-rel-perfect}
and find a pseudo-coherent $L \in D(\mathcal{O}_Y)$ which
``universally computes'' $Rf_*R\SheafHom(K, E)$ in the sense
described in that lemma. Unwinding the definitions, we obtain
for a point $y \in Y$ the equality
$$
\Ext^i_{\kappa(y)}(L \otimes_{\mathcal{O}_Y}^\mathbf{L} \kappa(y),
\kappa(y)) = \Ext^i_{\mathcal{O}_{X_y}}(K_y, E_y)
$$
We conclude that
$$
H^i(L \otimes_{\mathcal{O}_Y}^\mathbf{L} \kappa(y)) = 0
\text{ for } i > 0 \Leftrightarrow
\Ext^i_{\mathcal{O}_{X_y}}(K_y, E_y) = 0 \text{ for }i < 0.
$$
By Derived Categories of Schemes, Lemma \ref{perfect-lemma-jump-loci}
the set $W$ of $y \in Y$ where this happens defines an open of $Y$.
This open $W$ then satisfies the requirement in (1) for all morphisms
from spectra of fields, by the ``universality'' of $L$.

\medskip\noindent
Let's go back to $Y$ a general algebraic space.
Choose an \'etale covering $\{V_i \to Y\}$ by affine schemes $V_i$.
Then we see that the subset $W \subset |Y|$ pulls back to the corresponding
subset $W_i \subset |V_i|$ for $X_{V_i}$, $K_{V_i}$, $E_{V_i}$.
By the previous paragraph we find that $W_i$ is open, hence $W$ is open.
This proves (1) in general. Moreover, parts (2) and (3) are entirely formulated
in terms of the category $\textit{Spaces}/W$ and the restrictions
$X_W$, $K_W$, $E_W$. This reduces us to the case $W = Y$.

\medskip\noindent
Assume $W = Y$. We claim that for any algebraic space $V$ over $Y$
we have $Rf_{V, *}R\SheafHom(K_V, E_V)$ has vanishing cohomology
sheaves in degrees $< 0$. This will prove (2) because
$$
\Ext^i_{\mathcal{O}_{X_V}}(K_V, E_V) =
H^i(X_V, R\SheafHom(K_V, E_V)) = 
H^i(V, Rf_{V, *}R\SheafHom(K_V, E_V))
$$
by Cohomology on Sites, Lemmas
\ref{sites-cohomology-lemma-section-RHom-over-U} and
\ref{sites-cohomology-lemma-Leray-unbounded}
and the vanishing of the cohomology sheaves implies the
cohomology group $H^i$ is zero for $i < 0$ by
Derived Categories, Lemma \ref{derived-lemma-negative-vanishing}.

\medskip\noindent
To prove the claim, we may work \'etale locally on $V$.
In particular, we may assume $Y$ is affine and $W = Y$.
Let $L \in D(\mathcal{O}_Y)$ be as in the second paragraph of the proof.
For an algebraic space $V$ over $Y$ denote $L_V$ the derived pullback of
$L$ to $V$. (An important feature we will use is that $L$ ``works'' for all
algebraic spaces $V$ over $Y$ and not just affine $V$.)
As $W = Y$ we have $H^i(L) = 0$ for $i > 0$
(use More on Algebra, Lemma
\ref{more-algebra-lemma-lift-pseudo-coherent-from-residue-field}
to go from fibres to stalks). Hence $H^i(L_V) = 0$ for $i > 0$.
The property defining $L$ is that
$$
Rf_{V, *}R\SheafHom(K_V, E_V) = R\SheafHom(L_V, \mathcal{O}_V)
$$
Since $L_V$ sits in degrees $\leq 0$, we conclude that
$R\SheafHom(L_V, \mathcal{O}_V)$ sits in degrees $\geq 0$
thereby proving the claim. This finishes the proof of (2).

\medskip\noindent
Assume $W = Y$ but make no assumptions on the algebraic space $Y$.
Since we have (2), we see from
Simplicial Spaces, Lemma \ref{spaces-simplicial-lemma-fppf-neg-ext-zero-hom}
that the functor $F$ given by $F(V) = \Hom_{\mathcal{O}_{X_V}}(K_V, E_V)$
is a sheaf\footnote{To check the sheaf property
for a covering $\{V_i \to V\}_{i \in I}$ first consider the
{\v C}ech fppf hypercovering $a : V_\bullet \to V$ with
$V_n = \coprod_{i_0 \ldots i_n} V_{i_0} \times_V \ldots \times_V V_{i_n}$
and then set $U_\bullet = V_\bullet \times_{a, V} X_V$. Then
$U_\bullet \to X_V$ is an fppf hypercovering to which we may
apply Simplicial Spaces, Lemma
\ref{spaces-simplicial-lemma-fppf-neg-ext-zero-hom}.}
on $(\textit{Spaces}/Y)_{fppf}$. Thus to prove that $F$
is an algebraic space and that $F \to Y$ is affine and of
finite presentation, we may work \'etale locally on $Y$; see
Bootstrap, Lemma \ref{bootstrap-lemma-locally-algebraic-space-finite-type}
and
Morphisms of Spaces, Lemmas \ref{spaces-morphisms-lemma-affine-local} and
\ref{spaces-morphisms-lemma-finite-presentation-local}. We conclude
that it suffices to prove $F$ is an affine algebraic space of
finite presentation over $Y$ when $Y$ is an affine scheme. In this
case we go back to our pseudo-coherent complex $L \in D(\mathcal{O}_Y)$.
Since $H^i(L) = 0$ for $i > 0$, we can represent $L$ by a complex
of the form
$$
\ldots \to \mathcal{O}_Y^{\oplus m_1} \to \mathcal{O}_Y^{\oplus m_0}
\to 0 \to \ldots
$$
with the last term in degree $0$, see More on Algebra, Lemma
\ref{more-algebra-lemma-pseudo-coherent}. Combining the two displayed formulas
earlier in the proof we find that
$$
F(V) =
\Ker(
\Hom_V(\mathcal{O}_V^{\oplus m_0}, \mathcal{O}_V)
\to 
\Hom_V(\mathcal{O}_V^{\oplus m_1}, \mathcal{O}_V)
)
$$
In other words, there is a fibre product diagram
$$
\xymatrix{
F \ar[d] \ar[r] & Y \ar[d]^0 \\
\mathbf{A}_Y^{m_0} \ar[r] & \mathbf{A}_Y^{m_1}
}
$$
which proves what we want.
\end{proof}
